% !PRINTED 462
\weberpage{462}
If $h_1 = 2$, $h_2 = h_3 = h_4 = 0$, then $m = 12$. An
invariant curve $\Phi$ of the $12$th order must touch the curve $f$ at
the 24 points $P_7$, and if in $\Phi - a\Delta^2$ we determine the
constant $a$ suitably, it follows that this function
must be divisible by $f$. The quotient, as an invariant
of the $8$th order, can only be of the form $bf^2$, and accordingly
$\Phi$ is of the form $a\Delta^2 + bf^3$. In the same way we can conclude
that an invariant of the $18$th order must have the form
$a\Delta^3 + b\Delta f^3$.

\begin{center}
\gesperrt{All invariants of the $6$th, $12$th, $18$th
order are therefore rational functions of $\Delta$ and $f$.}
\end{center}

The next value that $m$ can have according to (1) is $m = 14$. In
this case $h_2 = 1$, while $h_1, h_3, h_4$ are equal to zero.
All invariants of the $14$th order therefore pass through the 56
points $P_3$, and these form the complete system of intersection points
of such an invariant curve with the fundamental curve. This holds also of
the invariant $C$. If we have a second invariant $C'$ of the $14$th
order, we can again, as above, determine the constant $a$ so
that $C' - aC$ is divisible by $f$. The quotient is
an invariant of the $10$th order, and from this it follows, since apart
from $f\Delta$ there is no invariant of the $10$th order:

\begin{center}
\gesperrt{Every invariant of the $14$th order is representable in the
form:}
\end{center}
\[
aC + bf^2\Delta,
\]
\begin{center}
\gesperrt{wherein $a, b$ are constants. Conversely, every expression of
this form is an invariant of the $14$th order.}
\end{center}

Next, $m$ can according to (1) have the value 20, namely for
$h_1 = h_2 = 1$. An invariant of the $20$th order must therefore pass
through the points $P_7$ and $P_3$, that is, through the 80
points of intersection of $f$ with $\Delta C$. From this we can conclude, just as
before, that an invariant of the $20$th degree is representable in
the form
\[
f(a\Delta^2 + bf^3) + cC\Delta.
\]
wherein $a, b, c$ are arbitrary constants. There is no
independent invariant of the $20$th order.

Let us now assume that $K'$ and $K$ are two invariants
of the $21$st order; both must by (1) pass through the 84 points
$\overline{P}_2$, and consequently one can determine $a$ so that
$K' - aK$ becomes divisible by $f$. The quotient would be an invariant
%CONT>>
